\section*{Chapter \ref{ch:ml:online-learning-and-drift} --- Online Learning and Drift}

\begin{solution}{exo:ml:online-learning-and-drift:1}
Memories $1/(1-\lambda)$: 100, 500 and 2\,000 observations. Relative weights $\lambda^{500}$: $0.99^{500} = 0.0066$,
$0.998^{500} = 0.37$, $0.9995^{500} = 0.78$.
\end{solution}

\begin{solution}{exo:ml:online-learning-and-drift:2}
$20\times1\,000/252 = 79.4$. The chapter's 74, 78 and 91 are within the Poisson noise of that count ($\pm9$), which is
what a calibration on another no-change stream can promise.
\end{solution}

\begin{solution}{exo:ml:online-learning-and-drift:3}
The statistic adds $\mu_0 - v - k$ each step, on average $0 + 0.38 - 0.1 = 0.28$; from zero it needs $9/0.28\approx32$
steps. The measured median is 17 because the statistic does not start at zero: under no change a reflected walk with
drift $-0.1$ and unit variance sits on average at about $1/(2\times0.1) = 5$, so at the change it is typically a few units
up already, and noise makes the first passage earlier still.
\end{solution}

\begin{solution}{exo:ml:online-learning-and-drift:4}
The error of a forgetting estimator has two parts: estimation noise, which falls as the memory $M$ grows (about $p/M$ in
R-squared units), and staleness, which grows with the share of the memory spent in past regimes (about $M/D$ times the
signal). Minimising $p/M + cM/D$ gives $M\propto\sqrt{pD/c}$: the best memory grows like the square root of the regime
length, not in proportion to it. The chapter's optimum, 200 observations at $D = 100$ and 500 at $D = 2\,000$, fits that
slow growth; at $D = 10\,000$ the stream holds only two regimes and forgetting nothing wins.
\end{solution}

\begin{solution}{exo:ml:online-learning-and-drift:5}
It was calibrated and timed on sign reversals, where the gains fall to about $-0.38$ standard deviations; when the
coefficients are redrawn at random the old model's gains fall only to about zero, a change of half the size, and a
threshold of 9 with a tolerance of 0.1 needs a long run of them. Recalibrate the tolerance for the change that matters
(fading to zero), watch a statistic that responds to it (the rolling correlation of forecast and outcome), and keep a
calendar refit as a floor under the detector.
\end{solution}

\begin{solution}{exo:ml:online-learning-and-drift:6}
The threshold was tuned on the two changes it is then judged on: its delay is measured in-sample and its false-alarm
rate is not measured at all. Two events cannot calibrate anything. Set the threshold on no-change data to a chosen
false-alarm rate, then measure the delay on changes the threshold never saw (planted ones, or later history).
\end{solution}

\begin{solution}{exo:ml:online-learning-and-drift:7}
With regimes of 500 steps on average: never refitted $-3.68\%$; every 250 steps $-0.60\%$ (75 refits); on a Page--Hinkley
alarm $-8.24\%$ (45 refits); \omterm{def:ml:online-learning-and-drift:rls}{recursive least squares} with $\lambda = 0.995$ 2.24\%; the truth 5.35\%. The calendar schedule's
500-observation window now spans a regime change most of the time, so its fits are stale; the detector alarms far more
often and each time refits on the 100 observations since the alarm, five coefficients from a sample too small for a signal
worth 5\% of the variance, and those noisy fits lose more than a stale one. Only the estimator whose memory suits the
regimes, about 200 observations, stays positive.
\end{solution}

\begin{solution}{exo:ml:online-learning-and-drift:8}
With $\lambda = 1$ the update of $P$ is the Sherman--Morrison formula for $P_t^{-1} = P_{t-1}^{-1} + x_tx_t^\top$, so
$P_t^{-1} = I/\delta + \sum_{s\le t}x_sx_s^\top$. By induction, $P_t^{-1}\hat\beta_t = \sum_{s\le t}x_sy_s$: it holds at
$t = 0$ with $\hat\beta_0 = 0$, and $P_t^{-1}\hat\beta_t = P_t^{-1}\hat\beta_{t-1} + P_t^{-1}k_t(y_t - x_t^\top\hat\beta_{t-1})$
with $P_t^{-1}k_t = x_t$ gives $P_{t-1}^{-1}\hat\beta_{t-1} + x_ty_t$. Hence $\hat\beta_t = (X^\top X + I/\delta)^{-1}X^\top y$,
the ridge estimator with penalty $1/\delta$, which tends to ordinary least squares as $\delta\to\infty$.
\end{solution}

\begin{solution}{pb:ml:online-learning-and-drift:1}
\textbf{Part I.}
\begin{enumerate}
\item 4.9\% to 5.4\% depending on the stream. With regimes of 100 steps the model must estimate five coefficients from
about a hundred observations before they change; its estimation noise is as large as the signal.
\item $\lambda = 0.995$ (memory 200) for $D = 100$ and 500, scoring 0.75\% and 2.44\%; $\lambda = 0.998$ (500) for
$D = 2\,000$, 3.71\%; $\lambda = 1$ for $D = 10\,000$, 5.21\% against the truth's 5.29\%.
\item 0.57\%, 2.05\%, 3.69\% and 4.84\%. A fixed \omterm{def:ml:trees-and-boosting:boosting}{learning rate} weights old observations geometrically, like a \omterm{def:ml:online-learning-and-drift:rls}{forgetting
factor}, but along one gradient direction at a time; it is a cheaper and noisier version of RLS.
\item That the true coefficients follow a random walk, with a variance per step set by $\lambda$: choosing $\lambda$ is
choosing how fast you believe the truth moves.
\end{enumerate}
\textbf{Part II.}
\begin{enumerate}[resume]
\item The gain $\hat y_ty_t$, standardised on the warm-up. It is what the desk earns from the model, and its mean changes
sign when the relation reverses; the squared error is dominated by the noise, which does not change.
\item $h = 9$ for CUSUM and Page--Hinkley (allowance and tolerance 0.1), $\delta = 0.0032$ for ADWIN.
\item CUSUM 20.8 (median 17), Page--Hinkley 27.4 (18.5), ADWIN 42.5 (26.5) steps.
\item 74, 78 and 91 over 20\,000 steps, against 79 expected: as calibrated. Whether one a year is too many is the desk's
cost question (Part IV).
\end{enumerate}
\textbf{Part III.}
\begin{enumerate}[resume]
\item Never $-7.81\%$; every 250 steps 2.62\%; on a Page--Hinkley alarm $-2.67\%$ (two refits); RLS with $\lambda = 0.998$
3.24\%; the truth 4.62\%.
\item Its coefficients belong to the first regime; in the other regimes they are unrelated or opposite to the truth, and a
forecast that is wrong in direction scores below the zero forecast.
\item It was tuned to reversals and the stream redraws; see \cref{exo:ml:online-learning-and-drift:5}.
\item It refits every step with the memory that suits the regimes, with no threshold to cross and no sample thrown away.
\end{enumerate}
\textbf{Part IV.}
\begin{enumerate}[resume]
\item \emph{The half-life of a model.} The \omterm{def:ml:online-learning-and-drift:rls}{forgetting factor} that minimises the prequential loss grows with the regime
length: $\lambda = 0.995$ at $D = 100$ and 500, 0.998 at 2\,000, 1 at 10\,000. At one false alarm per 252 steps, the median
delay to a sign reversal is 17 steps for CUSUM, 18.5 for Page--Hinkley and 26.5 for ADWIN.
\item An online estimator where the model is linear or has a linear last layer, with the memory chosen on history; a
calendar refit of the rest; a detector on the gains to trigger an early refit and a human review.
\item Scan the \omterm{def:ml:online-learning-and-drift:rls}{forgetting factor} or the rolling window on history and read the regime length off the optimum
(\cref{exo:ml:online-learning-and-drift:4}); or count the breaks Book~7's tests find.
\item A false alarm costs a refit (compute, validation, a new model's turnover) and trust; a missed change costs the
losses of a wrong model until the next refit, which for a reversed signal is the whole position's expected loss.
\item Book~7 measured a signal's decay after the fact; the \omterm{def:ml:online-learning-and-drift:rls}{forgetting factor} turns that half-life into the memory of the
estimator that trades the signal.
\item The time, the detector and its threshold, the statistic's path, the model version, and what was done: refit,
override, or dismissed, with a reason.
\item As a change in the features' distribution with the gains unchanged; detectors on the features (chapter~27) see it,
detectors on the gains do not until it hurts.
\item As fast as the truth moves and no faster, which you measure on history and check on the gains.
\end{enumerate}
\end{solution}

\begin{solution}{iq:ml:online-learning-and-drift:1}
\omterm{def:ml:online-learning-and-drift:drift}{Concept drift} changes $\P(y\mid x)$: a signal crowds out, or a fee change reverses the sign of an order-flow feature's
effect. \omterm{def:ml:online-learning-and-drift:drift}{Covariate shift} changes $\P(x)$ only: a volatility regime moves features into ranges the model rarely saw.
\iqlookfor{the two distributions and a correct example of each.}
\end{solution}

\begin{solution}{iq:ml:online-learning-and-drift:2}
Minimise $\sum_s\lambda^{t-s}(y_s - x_s^\top b)^2$; the normal equations $R_tb = r_t$ with $R_t = \lambda R_{t-1} + x_tx_t^\top$
and $r_t = \lambda r_{t-1} + x_ty_t$; Sherman--Morrison on $P_t = R_t^{-1}$ gives the gain, the coefficient update and the
covariance update of \cref{def:ml:online-learning-and-drift:rls}. $\lambda$ sets the memory, about $1/(1-\lambda)$
observations.
\iqlookfor{the weighted objective, the rank-one update, and the memory interpretation.}
\end{solution}

\begin{solution}{iq:ml:online-learning-and-drift:3}
On data with no change, pick the threshold that gives the false-alarm rate the desk can act on; then plant the changes that
matter and measure the delay; report both. Check on held-out history that the rate holds.
\iqlookfor{calibration on no-change data, a delay measured on unseen changes, and the trade-off between them.}
\end{solution}

\begin{solution}{iq:ml:online-learning-and-drift:4}
Score each schedule prequentially on history, costs of refitting included; estimate how fast the relation moves (the
forgetting-factor or window scan); weigh the gain against the operational risk of more frequent deployments; keep the
calendar as a floor and add a detector.
\iqlookfor{a prequential comparison, costs, and a drift estimate.}
\end{solution}

\begin{solution}{iq:ml:online-learning-and-drift:5}
It accumulates the deviations of the value below its running mean, less a tolerance, and alarms when the sum rises a
threshold above its running minimum. Weaknesses: the running mean includes all history since the last reset, so after a
long calm period it reacts slowly; it detects one direction and one size of change well, the one its tolerance and
threshold were tuned to.
\iqlookfor{the statistic and a real weakness (memory of the running mean, or tuning to one change).}
\end{solution}

\begin{solution}{iq:ml:online-learning-and-drift:6}
See \cref{prop:ml:online:kalman}: inflating the state covariance by $1/\lambda$ in the prediction step reproduces RLS. So
$\lambda$ encodes how fast the coefficients move relative to the noise; choose it by scoring on history, or estimate the
state noise by maximum likelihood and read off the implied $\lambda$.
\iqlookfor{the prediction step matched to the forgetting and the implication for choosing $\lambda$.}
\end{solution}
